\chapter{Verification Methodology}
\label{ch:verif-method}

% DRAFT BRIEF: explain HOW the design was verified before Ch.7 shows WHAT exists.
% Format = pseudocode + UVM topology diagram; avoid raw UVM source. Describe scoreboard/SVA
% from source --- never use the stale "minimal" label.
% Sources: verification specs 00/03/04/05, i3c_scoreboard.sv, UVM 1.2.

\section{Verification goals and phase split}
\brief{Closure goals: functional correctness of SDR write/read/immediate + \iic{}-FM + ENTDAA
+ ENEC/DISEC, scoreboard self-checking, SVA silent. Phase-1 = directed vseqs + scoreboard +
SVA; Phase-2 = constrained-random + functional coverage + broader CCC/error injection (future work).}

\tabph{Verification goals vs evidence}{spec 00}

\section{Directed-vs-constrained-random rationale}
\brief{Justify directed-first: protocol bring-up needs deterministic, debuggable stimulus to
pin each FSM path before randomization adds value. Constrained-random + coverage deferred to
Phase 2.}

\figph{Directed-vs-random rationale}{coverage-vs-effort tradeoff; where directed stops and CRV begins}
  {conceptual}{TikZ}

\section{Why UVM 1.2 + Xcelium}
\brief{Justify the toolchain explicitly (do not merely state it): UVM 1.2 (CDNS-1.2) for the
standard component/factory/TLM model; Xcelium (\texttt{xrun}) + SimVision for sim/debug. Note
the deviation from the supervisor outline's ModelSim/Verilator/GTKWave (stale). Justify a
custom register agent over \texttt{uvm_reg} given the simple addr/wen/ren bus.}

\tabph{Tool/methodology choices + rationale}{this work}

\section{Layered test stack}
\brief{test, virtual sequence, \{reg sequences, I3C device sequences\}, DUT --- with the
register agent driving the 32-bit bus and the I3C device-mode agent responding on the bus.
Pseudocode for a representative vseq.}

\figph{Layered test stack}{test to vseq to reg/i3c seqs to agents to DUT layering}{conceptual}{TikZ}

\section{TLM analysis path and scoreboard strategy}
\brief{Describe (from \texttt{i3c_scoreboard.sv}) the analysis path monitor to scoreboard and
the cross-checks: read/write data, T-bit parity, RX/TX byte packing, CCC, DAT entries,
SW-reset, response-priority, and end-of-test residue. Emphasise it is a substantive
scoreboard, not a stub.}

\figph{TLM path (monitor to scoreboard)}{monitor analysis ports to scoreboard checkers}
  {i3c_scoreboard.sv}{TikZ}

\section{SVA checkers}
\brief{Five SVA files in \texttt{i3c_core/sva/}: four bound (\texttt{flow_active_sva},
\texttt{csr_registers_sva}, \texttt{sync_fifo_sva}, \texttt{i3c_controller_top_sva}) +
\texttt{tb_pad_model_sva} instantiated. Summarise the properties each asserts; full binding
map in \cref{ch:uvm-env}.}

\section{Coverage strategy}
\brief{Phase-1 deferral stated honestly; Phase-2 blueprint: covergroups for CCC$\times$direction,
cmd-attr$\times$direction$\times$length, T-bit$\times$r/w, err_status, CSR addr$\times$access,
FIFO occupancy, DAT index$\times$is_i2c. If a \texttt{COV=1} run exists, point forward to
\cref{ch:results}.}

\section{Limitations / Phase-2 gaps}
\brief{Honest gaps: CRV + functional coverage not yet closed; CCC suite has 1 ad-hoc vseq;
multi-device + error-injection breadth limited; vseq suite is a WIP snapshot toward the
$\sim$105-case testplan. Frame all as future work.}
